% Author:
% Date: January 21, 2026

\subsection{Question Generation}
Based on Knowledge graph and synthetic patients, the Harness generates the following:

\subsubsection{Multiple Choice Question Generation}
- Topic selection
- Generate question Stem and right answer based on the graph
- Select distractors from the graph
- Otional: use LLM to form the question wording, otherwise use hard-coded templates.

\subsubsection{Vignette Based Question Generation}
\begin{verbatim}
llm.simulations.vignettes.VignetteGenerator
\end{verbatim}
Inputs:
Package name (HIV, immunization etc)
LLM model: a model to call to generate the vignette
Patient list: a list of patients to create vignettes. One question for each patient. If no list is provided will generate random patients.

Loops over a patient list to generate a question from one of the three topics selected randomly:
Differential diagnosis, initial management, or next diagnostic step
Output:
A list of dictionaries, one dict per question.

\subsubsection{Multi-Turn Conversation Generation}
- Simulate patient background
- LLM receives the patient background and is asked to impersonate patient. Behavior changes with patient age.
- An LLM health worker agent engages with the patient LLM.
- Three auxiliary LLMs work on the background handling: conversation closure (to avoid HW and patient for talking indefinitely), and lab results checker and generator (providing realism to the conversation).

Two methods:
\begin{verbatim}
llm.simulations.mt_conversations.make_conversation
\end{verbatim}

Base method for generating a conversation. Each LLM has to be provided as input.

\begin{verbatim}
llm.simulations.mt_conversations. simulate_conversation_default
\end{verbatim}

Instead, only pass the LLM to play the part of the health worker. All the other components are defaulted.
